% R911: gerado; alterar dados e gerador, não esta prancha.
\clearpage
\refstepcounter{figure}\label{fig:AR01-structure}
\noindent{\large\sffamily\bfseries\color{navy}AR01 · Diagrama de arquitetura}\par\smallskip
\noindent\textbf{Interfaces, configuração e entrada no Core}\par
\noindent\textit{Como o pedido chega ao Core por CLI, painel ou MCP?}\par\medskip
\begin{center}

\tikzset{
 atlasnode/.style={draw=steel,fill=blueBg,rounded corners=2mm,align=center,
  text width=4cm,inner sep=5pt,font=\fontsize{9.5}{11.8}\selectfont,text=textmain},
 atlasroot/.style={atlasnode,fill=navy,text=white,text width=2.8cm,font=\fontsize{10}{12}\selectfont\bfseries},
 atlasdecision/.style={diamond,aspect=2.1,draw=amber,fill=amberBg,align=center,
  text width=3cm,inner sep=2pt,font=\fontsize{9.5}{11.5}\selectfont},
 atlascall/.style={-Stealth,draw=steel,thick},
 atlasevent/.style={-Stealth,draw=violet,dashed,thick},
 atlasdata/.style={-Stealth,draw=teal,densely dotted,thick},
 atlasstructure/.style={draw=slate,thick},
 atlasboundary/.style={-Stealth,draw=amber,dash dot,thick},
 atlasedgelabel/.style={fill=papel,inner sep=1.3pt,font=\fontsize{9}{10}\selectfont}
}

\begin{tikzpicture}
\node[atlasnode] (c0) at (-4.3,0) {\textbf{C1: CLI de uso direto}\\Interpreta argumentos e seleciona operação};
\node[atlasnode] (c1) at (4.3,0) {\textbf{C2: Painel Streamlit}\\Obtém a instância usada pelo painel};
\node[atlasnode] (c2) at (-4.3,-3.0) {\textbf{C3: Ferramenta MCP de execução}\\Expõe execução limitada e cancelável};
\node[atlasnode] (c3) at (4.3,-3.0) {\textbf{C4: Contrato primário marceloclaro}\\Reúne serviços na instância do processo};
\node[atlasnode] (c4) at (-4.3,-6.0) {\textbf{C5: Coordenador autônomo}\\Conduz tarefa com limites e preflight};
\node[atlasnode] (c5) at (4.3,-6.0) {\textbf{C6: Configuração OpenCode}\\Compila agentes, comandos e MCPs};
\draw[atlascall] (c0.east) to[out=0,in=180] node[atlasedgelabel,pos=.25]{E1} (c3.west);
\draw[atlascall] (c1.south) -- node[atlasedgelabel,pos=.25]{E2} (c3.north);
\draw[atlascall] (c2.south) -- node[atlasedgelabel,pos=.25]{E3} (c4.north);
\draw[atlascall] (c4.east) to[out=0,in=180] node[atlasedgelabel,pos=.25]{E4} (c3.west);
\draw[atlasstructure] (c5.west) to[out=180,in=0] node[atlasedgelabel,pos=.25]{E5} (c0.east);
\draw[atlasstructure] (c5.west) to[out=180,in=0] node[atlasedgelabel,pos=.25]{E6} (c2.east);
\end{tikzpicture}
\end{center}
{\small\textbf{Figura \thefigure.} Diagrama de arquitetura --- Interfaces, configuração e entrada no Core}\par\smallskip
{\footnotesize\textbf{Chave das relações.} E1: C1 para C4 --- instancia no processo CLI (chamada); E2: C2 para C4 --- instância do processo web (chamada); E3: C3 para C5 --- solicita execução (chamada); E4: C5 para C4 --- delega quando pronto (chamada); E5: C6 para C1 --- configuração declarada (estrutura); E6: C6 para C3 --- servidores declarados (estrutura).}\par\smallskip
\noindent\textbf{Como funciona e por quê.} Localize as seis responsabilidades C1--C6. Os identificadores E1--E6 permitem conferir o significado de cada conexão na chave das relações.\par
\noindent\textbf{Justificativa.} Separar componentes e tipos de relação ajuda a localizar onde uma dependência, um evento ou uma fronteira foi declarado.\par
\noindent\textbf{Limite.} Não há uma instância global distribuída entre CLI, painel e servidores MCP. Interface aberta e comando disponível não comprovam autenticação, inferência ou entrega. A figura reúne rotas alternativas e contratos; não representa um traço executado.\par\smallskip
\fontecore{marceloclaro/cli.py}{main}
\fontecore{webapp/app.py}{get_orchestrator}
